\subsection{Learning Regimes and Restoration Systems}
\label{sec:learning}
\subsubsection{Paired, Self-Supervised, and Mixed Supervision}
The availability of training data determines how a restoration model connects its generative prior to degraded observations. Paired learning directly supervises this connection using degraded and clean videos. Because accurately aligned real pairs are difficult to acquire, synthetic degradation remains an important source of training data. It provides scalable supervision and precise control over corruption severity, but its effectiveness depends on how closely the simulated process represents real acquisition. A diffusion model trained on a narrow degradation distribution may reconstruct familiar inputs accurately while relying excessively on its prior when the observation falls outside that distribution.

For video, the degradation distribution has a temporal dimension. Independently sampling a corruption for each frame produces different statistics from persistent sensor noise or the periodic structure of a video codec. Mixed-degradation training further involves interactions among corruption types, strengths, and processing order. Noise followed by compression, for example, produces a different observation from compression followed by noise. Consequently, generalization concerns both unfamiliar individual degradations and unfamiliar combinations or temporal patterns of otherwise familiar degradations.

Self-supervision can complement paired learning at different stages of the restoration pipeline. SCST learns degradation-insensitive control features through a self-supervised objective within a broader paired VSR pipeline~\cite{scst}. UnDIVE combines unlabeled diffusion pretraining, paired spatial enhancement, and an unsupervised temporal constraint~\cite{undive}. These methods illustrate a hybrid learning regime in which prior learning, spatial reconstruction, and temporal regularization draw on different forms of supervision. Describing supervision at the component level therefore provides a more informative account of data requirements than assigning a single supervised or unsupervised label to the entire system.

InstantViR~\cite{instantvir} distills a pretrained video prior into a feed-forward inverse solver using degraded measurements and a known forward operator, without external clean/degraded training pairs. A measurement-likelihood term constrains fidelity while the frozen teacher regularizes the solution distribution. The solver uses bidirectional attention within each temporal block and causal attention with cached keys and values across blocks; its streaming claim therefore includes block buffering. LeanVAE reduces reconstruction overhead, with teacher-space regularization compensating for the change of latent representation during training. This connects generative distillation to blockwise streaming reconstruction.

\subsubsection{Zero-Shot Transfer and Test-Time Adaptation}
Zero-shot restoration encompasses several forms of transfer: applying an image restoration model to video, using a general generative prior with a new degradation operator, and restoring an unseen domain without paired target-domain training. ZVRD and DiffIR2VR-Zero exemplify image-to-video transfer~\cite{zvrd,diffir2vr}. Their relevance lies in reusing spatial restoration knowledge while introducing temporal coupling at inference. Zero-shot transfer does not necessarily eliminate auxiliary information: optical flow, pretrained restoration modules, estimated degradation parameters, and test-time optimization can still play substantial roles.

Test-time adaptation addresses distribution shift by updating model parameters from the observed test video. It differs from optimizing latent variables or initial noise while keeping the model fixed. Diffusion test-time adaptation for adverse weather combines a diffusion formulation with temporal calibration~\cite{difftta}. Such adaptation offers a way to exploit sequence-specific information, although its effectiveness depends on the amount and reliability of the available observations. Short sequences, abrupt scene changes, and severely corrupted frames can weaken the adaptation signal, while access to future frames changes whether a method is suitable for online processing.

These distinctions also clarify the computational implications of training-free inference. A frozen model may still require many gradient evaluations to optimize its input noise or latent representation. Conversely, a small parameter update can provide inexpensive sequence-specific calibration. The practical difference between these approaches thus involves the variables being optimized, the available temporal context, and the additional inference cost, rather than the presence or absence of offline training alone.

Temporal-Consistent Video Restoration~\cite{tcvr} instead treats the full reverse diffusion map as a differentiable generator and optimizes a shared seed and frame-dependent residuals under a MAP objective. Semantic structure in seed space and progressive optical-flow warping address consistency at different levels. Its four reverse steps describe the inner generator; the reported reconstruction still involves thousands of outer optimization iterations. The short inner chain makes differentiation through the complete generator tractable.

Domain-adaptive Video Deblurring via Test-time Blurring~\cite{dadeblur} uses diffusion to synthesize degradation rather than directly restore the test video. Relatively sharp patches are extracted from the target domain, and a pre-existing ID-Blau blur generator produces domain-matched pseudo training pairs that adapt a separate deblurrer. The resulting DADeblur framework uses generative modeling to improve the training distribution of a conventional restorer. Diffusion generates adaptation data at test time; after adaptation, the separate deblurring network restores frames without a reverse-diffusion chain for each output.

\subsubsection{Single-Task, Multi-Task, and All-in-One Models}
Single-task restoration allows the observation encoder, degradation model, and training objective to specialize in one problem. Multi-task learning shares representations across problems, while all-in-one restoration aims to handle multiple degradations with a common model, often through degradation prompts or shared conditional features. This distinction is particularly relevant when one architecture is trained into several task-specific checkpoints: architectural reuse provides less operational unification than a single checkpoint that handles all supported tasks. TDM and generic restoration models such as SeedVR illustrate the move toward broader task coverage~\cite{tdm,seedvr}; UniFlowRestore pursues a related objective through flow matching~\cite{uniflowrestore}.

Task diversity introduces interactions that cannot be inferred from performance on isolated degradations. Denoising may remove detail useful for subsequent deblurring, whereas deblurring may amplify residual noise. A model that handles each task independently therefore does not necessarily handle their composition. Open-degradation restoration extends this difficulty to corruption types and combinations absent from training. In this setting, identifying the degradation and estimating the reliability of the observation become closely connected to controlling the strength of generative correction.

Video adds a further interaction between task identification and temporal consistency. If degradation representations or routing decisions fluctuate across adjacent frames, restoration behavior can change even when the underlying scene remains stable. Temporally coherent conditioning helps avoid such discontinuities, but excessive persistence can delay adaptation to genuine changes in illumination, acquisition conditions, or scene content. Unified video restoration consequently requires both shared spatial restoration capability and appropriate temporal behavior in its degradation representation.

\subsubsection{Agentic and Routed Restoration}
Agentic restoration extends unification from a single model to a collection of tools. MoA-VR and VQ-Jarvis investigate systems that identify degradations, select or compose restoration operations, and assess output quality~\cite{moavr,vqjarvis}. Such systems can integrate diffusion models with regression networks, interpolation, compression, and quality estimation. Their relevance to diffusion-based restoration is therefore at the level of orchestration; diffusion need not be the mechanism of every component.

The effectiveness of orchestration depends jointly on the available specialists and the policy that selects them. Comparisons with a fixed operator or a fixed processing sequence help distinguish the contribution of adaptive routing from the strength of the tool library. An oracle selection provides an additional reference when the candidate outputs can be evaluated reliably. Tool calls and cumulative processing cost are also central to this system paradigm because repeated assessment and refinement can substantially increase latency.

Quality feedback creates a connection between restoration agents and the evaluation issues discussed in Section~\ref{sec:evaluation}. A judge that favors sharpness may encourage repeated enhancement and accumulate halos or unsupported texture; a judge that favors temporal smoothness may suppress genuine motion. Evidence preservation, stopping criteria, and recovery from harmful intermediate operations are therefore important research problems for agentic video restoration. They connect the choice of restoration tools to the reliability of the quality model that guides their use.
